\section{Store records in a circular byte sequence}
A ring buffer reuses a fixed allocation by returning to its first byte after its
last byte. The physical addresses still form an ordinary contiguous array. Circular
behavior comes from the program's index calculation. No memory address wraps
implicitly when an instruction advances a pointer beyond the allocation.

This logger stores a record of length \(L\) as one prefix byte containing \(L\),
followed by \(L\) payload bytes. Its encoded size is \(N=L+1\). There is no terminator,
alignment padding, or unused sentinel byte between records. A zero-length record
still occupies its prefix byte. Its payload is empty, but the record exists and
can be removed from a nonempty queue.

\Needspace{14\baselineskip}
\begin{longtable}{rll}
  \toprule Offset & Storage    & Meaning                                      \\\midrule
  0               & Word       & Head: index of the oldest record's prefix    \\
  4               & Word       & Tail: index where the next record will begin \\
  8               & Word       & Used: occupied bytes, including all prefixes \\
  12              & Word       & Capacity \(C\), one through 32 bytes         \\
  16              & Doubleword & Full address of the ring's first byte        \\\bottomrule
\end{longtable}
The first three words are mutable queue state. Capacity and buffer address describe
the allocation and remain unchanged. The provided physical ring arrays have 32
bytes, but a descriptor can expose a smaller initial span. Bytes beyond its capacity
are mapped padding, not part of that queue. All ring reads and writes must fit
within the exposed span. A pointer to the descriptor is distinct from the buffer
pointer stored inside it; load the latter with \texttt{ld} before addressing bytes.

The descriptor is a native in-memory object in MARS's default little-endian layout.
It is not a network packet or portable file format. Its fields are naturally aligned
by the supplied scaffold. Individual record bytes can cross the ring's physical end.
Read and write them using byte operations and explicitly wrapped indices.
